\subsubsection{The product order}
In each strand we compare monomials of the same bidegree $(n,d)$.
Write them as $x_{\boldsymbol a}=x_{a_1}\cdots x_{a_d}$ and
$x_{\boldsymbol b}=x_{b_1}\cdots x_{b_d}$, so that
$\sum_j a_j=\sum_j b_j=n$. On the odd primes our order is simply
\[
 x_{\boldsymbol a}\prec x_{\boldsymbol b}
 \quad\Longleftrightarrow\quad
 a_1\cdots a_d<b_1\cdots b_d.
\]
Unique factorization distinguishes all the monomials in the fiber.

The connection with \emph{majorization} explains why balanced pairs
give the largest products. List both partitions in decreasing order.
Then $\boldsymbol a$ majorizes $\boldsymbol b$ if
$\sum_{j=1}^s a_j\ge\sum_{j=1}^s b_j$ for each $1\le s<d$.
For distinct partitions of positive labels this implies
$a_1\cdots a_d<b_1\cdots b_d$, by strict concavity of $\log$ and
the majorization inequality \cite[Ch.~3]{MOA11}.
For $d=2$ the larger of the two parts determines majorization.
For $d=3$ the product also compares partitions that majorization
leaves incomparable. In $R_{31,3}$, for instance, the first two
partial sums for $(13,13,5)$ are $13,26$, whereas those for
$(17,7,7)$ are $17,24$. Their products are $845$ and $833$.

For an arbitrary bounded-below support $\A$, put
$c=\max\{0,2-\min\A\}$ and compare first by
$P_c(x^\alpha)=\prod_{a\in\A}(a+c)^{\alpha(a)}$.
The same majorization argument applies to $\log(a+c)$.
If two products agree, take the largest label $a_0$ at which the
exponents differ and put $x^\alpha\prec x^\beta$ if
$\alpha(a_0)<\beta(a_0)$. For example, $x_2x_6^2$ and $x_3^2x_8$
on $\A=\{2,3,6,8\}$ have the same sum, monomial degree and product.
The comparison by $\log P_c$ is a weight order with weights
$\log(a+c)$ \cite[Definition~1.2 and Section~6]{Vietri02}.

\begin{proposition}[Product monomial order]\label{moves:order}
The same comparison on all monomials is a monomial well-order on $R$.
With $c$ fixed, it restricts to each polynomial subring on a subset
of the variables.
\end{proposition}
\begin{proof}
The product comparison followed by the stated lexicographic tie-break
is a total order, with $1$ its least element. Products multiply, and
adding the same exponent vector preserves the largest differing label
and its exponent comparison. Hence multiplying by any monomial
preserves strict inequalities.

For $N\ge1$, only finitely many monomials satisfy $P_c\le N$.
Every factor satisfies $2\le a+c\le N$, and the number of factors
is at most $\log_2 N$. Every nonempty set of monomials therefore has a
least element, obtained by taking its least product score and then
the least of the finitely many ties. Restriction changes none of
these comparisons.
\end{proof}

\subsubsection{Minimal binary Goldbach moves}
For every represented sum $n$, let
\[
 \mathcal M_{n,2}=\{x_ax_{a'}:a,a'\in\A,\ a\le a',\ a+a'=n\}.
\]
Choose $b\le b'$ in $\A$ with $b+b'=n$ and $b'-b$ as small as
possible. We call $u_n=x_bx_{b'}$ the \emph{anchor} of this fiber.
If $n/2\in\A$, the anchor is the square $x_{n/2}^2$.

\begin{proposition}[Anchors and minimal generators]\label{moves:anchors}
The anchor $u_n$ is the unique largest monomial in $\mathcal M_{n,2}$
for the product order, independently of the shift $c$. The binomials
\[
 u_n-x_ax_{a'},\qquad
 x_ax_{a'}\in\mathcal M_{n,2}\setminus\{u_n\},
\]
as $n$ varies, form a minimal bihomogeneous generating set of $I$
over every field. If a quadratic fiber is nonempty in $R^{<m}$, its
anchor already belongs to $R^{<m}$. Consequently these chosen
generating sets are nested under variable truncation.
\end{proposition}
\begin{proof}
For $a+a'=n$ and $a\le a'$,
\[
 (a+c)(a'+c)=\frac{(n+2c)^2-(a'-a)^2}{4}.
\]
The sum and the gap determine the pair. Thus $u_n$ is the unique
product maximum for every $c$. For any other pair we have
$a<b\le b'<a'$: decreasing the gap at fixed sum raises the smaller
label and lowers the larger one. Hence if $a'\le m$, both $b$ and $b'$
are at most $m$. This proves the assertion about cutoffs.

For any two monomials $v,w\in\mathcal M_{n,2}$, their difference
is $(u_n-w)-(u_n-v)$. Within a fiber, the displayed binomials are
linearly independent because each nonanchor monomial appears in just one.
Different fibers have different Goldbach degrees. Since all generators
have monomial degree two and $\m I$ has monomial degree at least
three, their classes form a basis of $I/\m I$. This proves
minimality. The fixed anchors give the asserted nesting.
\end{proof}

\begin{definition}[Minimal binary Goldbach move]\label{moves:minimal}
A \emph{minimal binary Goldbach move} is a replacement
\[
 x^\alpha x_ax_{a'}\longrightarrow x^\alpha u_n,\qquad
 x_ax_{a'}\in\mathcal M_{n,2}\setminus\{u_n\},
\]
where $x^\alpha$ is the fixed monomial factor. The arrow points
from the nonanchor multiple to the anchor multiple, and hence
increases the product order. We also allow traversal in the opposite
direction when discussing undirected paths.
\end{definition}

Here ``minimal'' refers to the chosen minimal generating set of the
ideal. When we wish to specify a quadratic generator by the monomial
$v=x_ax_{a'}$ it replaces, we write $\left.e\right|_v$. Thus
$\partial(\left.e\right|_v)=u_n-v$ for $n=a+a'$.

\begin{proposition}[The vertices and edges of a strand]\label{moves:graph}
Choose these quadratic generators as the basis of $F_1$.
For every bidegree $(n,d)$,
the map $F_{1,n,d}\to F_{0,n,d}$ is the oriented incidence map
of the graph whose vertices are the monomials $x^\alpha$ of Goldbach degree
$n$ and monomial degree $d$
and whose edges are the minimal binary Goldbach moves in that fiber.
Its connected components are those of $\Gamma_{n,d}$.
\end{proposition}
\begin{proof}
Since $F_{0,n,d}=R_{n,d}$, its basis consists of the stated monomials.
The basis of $F_{1,n,d}$ consists of the multiples
$x^\alpha(\left.e\right|_v)$ of bidegree $(n,d)$. For $v=x_ax_{a'}$,
\[
 \partial\bigl(x^\alpha(\left.e\right|_v)\bigr)
     =x^\alpha u_{a+a'}-x^\alpha x_ax_{a'}.
\]
This is exactly the incidence boundary of the indicated edge.
Any elementary binary move between two multiples of pairs of sum
$a+a'$ with the same fixed factor $x^\alpha$ passes through $x^\alpha u_{a+a'}$
using at most two star edges. Thus the chosen edges connect the
vertices in each component of $\Gamma_{n,d}$.
\end{proof}

Along an arrow $x^\alpha x_ax_{a'}\to x^\alpha u_n$, the pair $(a',a)$ majorizes
$(b',b)$. Appending the labels of $x^\alpha$ to both pairs
preserves majorization: the sum of the $j$ largest labels is the
maximum over all ways to choose $j$ labels from the pair and from $x^\alpha$.
Thus the source partition majorizes the target along every directed path.
Every arrow strictly increases the product order, so the move graph
has no directed cycles. Its reachability relation is a partial order,
but the graph need not be its Hasse diagram: a chosen edge may also
be traversed by a longer directed path. Nor do these quadratic
generators have to be a Gr\"obner basis. The product order selects
and orients the edges; choosing higher syzygies and realizing their
boundaries by cells is a further problem.
Majorization alone does not create an edge: $\Gamma_{n,d}$ allows all
binary replacements, while the chosen graph retains only anchor moves.

\begin{proposition}[Maximal representatives]\label{moves:maximal-representatives}
Every connected component $C$ of $\Gamma_{n,d}$ has a unique greatest
monomial $m_C$ in the product order, called its \emph{component maximum}.
The classes of these monomials form a basis of $G_{n,d}$.
From any vertex, successive product-increasing minimal binary moves
terminate at a \emph{one-move local maximum}: a monomial from which
every nontrivial binary move strictly decreases $P_c$.
\end{proposition}
\begin{proof}
The fiber is finite, so each component has a greatest element in the
total product order, and a strictly increasing sequence of moves
terminates. At a terminal vertex every pair of factors is its sum's
anchor; otherwise replacing that pair by its anchor would increase
the product. Every nontrivial replacement of an anchor decreases
$P_c$. The basis assertion is Proposition~\ref{setup:components}.
\end{proof}

\begin{example}[Local maxima need not be unique]\label{moves:local-maxima}
On the odd primes, the minimal moves in bidegree $(21,3)$ include
\[
 x_5^2x_{11}\;\longleftrightarrow\;
 x_3x_7x_{11}\;\longleftrightarrow\;x_7^3,
 \qquad 275\;>\;231\;<\;343.
\]
The displayed numbers are the products of the labels. Both endpoints
are one-move local maxima: their quadratic factors are the anchors
$u_{10}=x_5^2$, $u_{16}=x_5x_{11}$, and $u_{14}=x_7^2$.
The component maximum is $x_7^3$, since the product of three positive
labels with sum $21$ is at most $7^3$. Thus reaching the component
maximum from a local maximum can require a product-decreasing move.
\end{example}

